\begin{claim}{C.81}{Proved}{Exact metric, compact-support, and finite-configuration comparisons for discrete operators}
Let $O=\operatorname{int}\overline D$ be the actual
open disk, let $d$ be a compatible metric on its
closed disk, and let $Q\leq\operatorname{Homeo}(O)$
be closed, discrete, free and proper. Retain the
actual relative centralizer
\[
 \mathscr C_Q=
 \{h\in\operatorname{Homeo}(\overline D):
 h|_{\partial\overline D}=\id,
 \ hq=qh\text{ on }O\text{ for every }q\in Q\}
                                                               \tag{C.81.1}
\]
with compact-open convergence of maps and inverses.
Assume that an actual subgroup $G\leq\mathscr C_Q$
is transitive on $O$. Write $N=O/Q$ and retain
the original quotient covering $\pi:O\to N$.
No extension of the operators $q$ to the circle
is assumed. The following comparisons hold.

The metric
\[
 \begin{split}
 d_Q(x,y)&=\sup_{q\in Q}d(qx,qy),\\
 \rho(\pi x,\pi y)&=\inf_{q\in Q}d_Q(x,qy).
 \end{split}
                                                               \tag{C.81.2}
\]
gives the original topology on $O$ and its
quotient topology on $N$. The specified covering
is locally isometric for these metrics and each
fiber is uniformly discrete. Restriction embeds
$\mathscr C_Q$ as a closed subgroup of the uniform
$d_Q$-homeomorphisms commuting with $Q$, and is
an isometry for the sum of the uniform map and
inverse distances and the corresponding actual
closed-disk distances.

Every compact parameter family of equivariant
interior homeomorphisms and their inverses,
supported over one fixed compact subset of $N$,
extends jointly by the identity to the original
circle. In particular, every based sphere family
in $\mathscr C_Q$ whose descents are supported
in the interior of one compact connected
subsurface $N_0\subset N$ with nonempty boundary
has a based filling in $\mathscr C_Q$, retaining
that quotient support and every original label.
Here sphere dimension is any $n\geq1$.

For finite $A\subset O$ with distinct quotient
labels put
$\mathscr C_Q(A)=\{h\in\mathscr C_Q:h|_A=\id\}$.
For $x\notin QA$, the actual evaluation map is
a locally trivial principal bundle
\[
 \mathscr C_Q(A\cup\{x\})\longrightarrow
 \mathscr C_Q(A)\xrightarrow{\ h\mapsto h(x)\ }O\setminus QA.
                                                               \tag{C.81.3}
\]
Its source and fiber retain their original
closed-disk topologies. Every inclusion
$\mathscr C_Q(A)\hookrightarrow\mathscr C_Q$
induces an isomorphism on $\pi_n$ at $1$ for
every $n\geq1$.

When $N$ is noncompact, descent from the subgroup
of actual lifts of compactly supported quotient
isotopies is not continuous from its induced
$\mathscr C_Q$ topology to the Whitney topology
on the quotient homeomorphism group. The proof
gives an explicit convergent sequence and one
positive continuous majorant excluding all its
descents.

These are finite-configuration and fixed compact
quotient-support comparisons. They do not provide
a coherent infinite-configuration homotopy or
a controlled fragmentation of every original
family into compact quotient supports.
\end{claim}

\begin{proof}
\emph{The invariant metric and the actual open evaluation map.}
Let $L$ be the closure of $G$ in (C.81.1).
It is a Polish group, remains transitive on $O$,
and commutes with every $q$. Closedness of the
commutator equations uses only the interior:
if $h_j\to h$ with inverses, then $h_jx\to hx\in O$,
so continuity of $q$ gives $qh_jx\to qhx$.
The verified Polish orbit theorem in C.30 makes
$L\to O$, $g\mapsto gx$, open.

The first formula (C.81.2) is finite because
$d$ is bounded on the compact disk. The identity
term separates points, and the triangle inequality
holds term by term. Reindexing by right
multiplication proves $Q$-invariance. Also
$d(x,y)\leq d_Q(x,y)$.
For the reverse topology comparison, take
$x_j\to x$ in $O$. Openness of evaluation
supplies $g_j\to1$ in $L$ with $g_jx=x_j$:
choose a decreasing identity-neighborhood basis
$U_m$, let $m(j)\to\infty$ so that
$x_j\in U_{m(j)}x$, and choose one such $g_j$.
No section has been assumed. Commutation gives
\[
 d_Q(x_j,x)=\sup_{q\in Q}d(g_jqx,qx)
       \leq\sup_{z\in\overline D}d(g_jz,z)
       \longrightarrow0.
                                                               \tag{C.81.4}
\]
This proves equality of the two interior
topologies, and in particular separability
of the displayed invariant metric space.

\emph{The specified quotient is a metric covering.}
Properness, discreteness and freeness give an
open neighborhood $U$ of any $x$ whose distinct
$Q$-translates are disjoint. Choose $r>0$ with
$B_{d_Q}(x,r)\subset U$. Then
$d_Q(x,qx)\geq r$ for $q\ne1$, and invariance
makes the entire fiber $Qx$ uniformly $r$-discrete.

The formula for $\rho$ is independent of both
representatives by invariance and reindexing.
Symmetry and the triangle inequality follow
from those of $d_Q$. A zero value gives
$q_jy\to x$ for some $q_j$; the closed-orbit
property of the proper action then gives
$x\in Qy$. Thus it is a metric. The equality
\[
 \pi B_{d_Q}(x,\epsilon)=B_\rho(\pi x,\epsilon)
                                                               \tag{C.81.5}
\]
shows that it gives the quotient topology,
and $\rho(\pi x,\pi y)\leq d_Q(x,y)$.
For $y,z\in B_{d_Q}(x,r/5)$ and $q\ne1$,
\[
 d_Q(y,qz)\geq r-2r/5=3r/5,
                    \qquad d_Q(y,z)<2r/5.
                                                               \tag{C.81.6}
\]
Consequently $\rho(\pi y,\pi z)=d_Q(y,z)$.
The inverse image of $\pi B_{d_Q}(x,r/5)$
is precisely the disjoint union of the translates
of that ball, each mapped isometrically onto
this quotient neighborhood. These are all
three metric-cover conditions in
\cite[Definition~2.1]{C_YagasakiUniform2014},
proved here for the actual map $\pi$.
The separation constant may depend on the
fiber; no compact quotient is inferred.

\emph{The exact isometry of the relative source group.}
Let $\mathscr U_Q$ denote the group of uniform
homeomorphisms of $(O,d_Q)$ commuting with $Q$,
with the uniform topology on maps and inverses.
An $h\in\mathscr C_Q$ belongs to this group
after restriction. Indeed its modulus of
uniform continuity on the original compact
disk applies to every pair $qx,qy$; commutation
then bounds $d(hqx,hqy)$ uniformly in $q$.
This is uniform $d_Q$-continuity of $h$.
Apply the identical argument to $h^{-1}$.

For $h,k\in\mathscr C_Q$ the exact formula is
\[
 \begin{split}
 \sup_{x\in O}d_Q(hx,kx)
   &=\sup_{x\in O,\ q\in Q}d(hqx,kqx)\\
   &=\sup_{z\in\overline D}d(hz,kz).
 \end{split}                                                   \tag{C.81.7}
\]
The last equality uses $q=1$ and density of
$O$. Repeat it with both inverses. Thus
restriction is an isometry for the stated
sum metrics, and its inverse on its image
is continuous. The source metric is complete:
simultaneous uniform limits of maps and
inverses on the compact disk are continuous
inverse maps, and all boundary and commutator
equations are closed. Its isometric image
is therefore closed in $\mathscr U_Q$.

This does not give an extension for arbitrary
members of $\mathscr U_Q$. For the explicit
case $Q=\{1\}$ and the standard round disk,
$d_Q=d$ and a nonidentity rotation is a
uniform interior homeomorphism. Its boundary
value is the same rotation, so it has no
preimage under this particular restriction
map. The named relative image condition
must consequently be retained in any use
of a uniform-homeomorphism deformation.

\emph{Compact cores and one common original source bound.}
Fix $x_0\in O$. C.64's proved compact-covering
lemma for the actual open evaluation map
gives, for every compact $K\subset O$, a
compact family $S\subset L$ with $K\subset S(x_0)$.
One covers $K$ by finitely many evaluation
neighborhoods and takes the union of their
compact preimages. For compact $K,K'$ choose
such $S,T$, respectively. For $z=sx_0\in K$
and $z'=tx_0\in K'$, commutation gives
\[
 \begin{split}
       qz'&=tqx_0=ts^{-1}(qz),\\
       qK'&\subset TS^{-1}(qz)\quad(q\in Q,\ z\in K).
 \end{split}
                                                               \tag{C.81.8}
\]
Every member of the one compact family
$TS^{-1}$ is an original circle-fixed map.
Thus $z_j\in K$ and $q_jz_j\to y$ on the
original circle imply uniform convergence
of the whole sets $q_jK'$ to $y$. No boundary
extension or convergence of $q_j$ is used.

Now take a compact parameter family $h_p$ of
$Q$-equivariant interior maps and inverses,
identity outside $\pi^{-1}(K_0)$ for a fixed
compact $K_0\subset N$. Choose compact
$C\subset O$ with $\pi C=K_0$, using finitely
many local sections of the covering. Include
in $C'$ all of $C$, $h_p(C)$ and $h_p^{-1}(C)$.
Joint interior continuity and compactness make
$C'$ compact. Since $\pi^{-1}(K_0)=QC$,
(C.81.8) supplies one compact original
family $P=TS^{-1}\cup\{1\}$ with
\[
 \begin{gathered}
                h_p(x)\in P(x),\qquad h_p^{-1}(x)\in P(x)\\
                       (p\text{ any parameter},\ x\in O).
 \end{gathered}
                                                               \tag{C.81.9}
\]
Inside $QC$ use equivariance and (C.81.8);
outside it use the identity. Compact evaluation
of this boundary-fixed $P$ proves both joint
circle limits. The pasted maps and inverses
are inverse homeomorphisms of the actual disk,
continuous in the actual topology of (C.81.1).

\emph{A based sphere with fixed compact quotient support.}
Let $f:(S^n,*)\to(\mathscr C_Q,1)$, $n\geq1$,
have descents $b_\xi$ identity outside
$\operatorname{Int}N_0$, where $N_0$ is as
in the statement. Descent and its inverse
family are continuous by the open covering
quotient and its local sections. The restrictions
to $N_0$ lie in the identity component of
$\operatorname{Homeo}(N_0\mathrel{\mathrm{rel}}\partial N_0)$:
the sphere is path connected and its basepoint
is the identity.

That component is contractible by
\cite[Lemma~2.1(ii)]{C_Yagasaki2000}.
The nonempty entire boundary is fixed, so
all its compact-surface exceptions are excluded.
A contraction can be normalized to fix the
identity throughout by replacing its value
$C_r(h)$ by $C_r(h)C_r(1)^{-1}$, with the
parameter directed from $1$ to $h$.
Applied to $b_\xi$ and extended by the identity
off $N_0$, it gives a continuous family
\[
 B_{0,\xi}=1,\qquad B_{1,\xi}=b_\xi,
                         \qquad B_{r,*}=1
                                                               \tag{C.81.10}
\]
of homeomorphisms of $N$ supported in $N_0$.

Lift the homotopy $B_{r,\xi}\pi$ from its
specified initial lift $H_{0,\xi}=\id_O$.
The covering homotopy property gives a joint
lift on $[0,1]\times S^n\times O$; it does
not require the sphere parameter to be simply
connected. Lift the inverse homotopy as well.
The two composites lift the identity and
start at the identity, hence are the identity
by uniqueness of lifts. For each original
$q$, the maps $H_{r,\xi}q$ and $qH_{r,\xi}$
lift the same family with identical initial
value $q$. They agree. The basepoint family
in (C.81.10) lifts to the constant identity.

At $r=1$, both $H_{1,\xi}$ and $f_\xi$ lift
the same $b_\xi$. Their relative deck label
is locally constant in $\xi$ and equals $1$
at $*$. Connectedness of the sphere therefore
gives $H_{1,\xi}=f_\xi$ for every parameter.
This proves agreement with the original
family upstairs, rather than only after
descent. Points above $N\setminus\operatorname{Int}N_0$
have constant base tracks and initial identity
lifts, so they stay fixed. Formula (C.81.9)
applies to this compact parameter family and
its inverse. The resulting joint original-circle
extension is the asserted based filling in
$\mathscr C_Q$ with every deck label retained.

\emph{Continuous local sections for actual orbit evaluation.}
For finite $A$ as in (C.81.3), its $Q$-saturation
is closed and locally finite in $O$. Commutation
identifies $\mathscr C_Q(A)$ with the subgroup
fixing that saturation pointwise. Choose
$x\notin QA$ and a small quotient disk $U$
about $\pi x$, with compact closure disjoint
from $\pi A$, that is evenly covered.
A local coordinate bump translation moves
$\pi x$ to each sufficiently nearby $y$,
using a fixed compact subdisk as support.
One explicit local formula is
$z\mapsto z+\chi(z)v$, where $\chi=1$ near
the center, $\chi=0$ outside that subdisk,
and $|v|\operatorname{Lip}(\chi)<1/2$.
The two-sided Lipschitz estimate makes it
injective with continuously varying inverse;
properness and invariance of domain make
it onto. The parameter $av$, $0\leq a\leq1$,
gives its based supported isotopy.

Lift that isotopy from the identity, retaining
the sheet containing $x$. It fixes $QA$,
commutes with every $q$, and its endpoint
maps $x$ to the corresponding local lift
of $y$. The compact quotient-support proof
above gives its continuity with inverses
in the original closed-disk topology.
Denote this actual local section by $\sigma_y$,
now with $y$ denoting the upstairs point,
and choose $\sigma_x=1$.

The evaluation map is onto even from its
identity path component. Given any other
$x'\in O\setminus QA$, choose a path from
$x$ to $x'$ in that connected punctured plane.
Cover its compact projected path by finitely
many quotient coordinate disks avoiding $\pi A$.
Their short supported point-moving isotopies
compose to move the point along that path.
The upstairs path is the chosen one by
uniqueness of the original covering lift.
All supports lie in one compact quotient
set, so their lifted family extends by
(C.81.9) and fixes $QA$. Its endpoint sends
$x$ to $x'$ in the asserted original group.

Set $F=\mathscr C_Q(A\cup\{x\})$. The local
bundle chart and its inverse are exactly
\[
 \begin{split}
 h&\longmapsto\bigl(h(x),\sigma_{h(x)}^{-1}h\bigr),\\
 (y,k)&\longmapsto\sigma_y k\qquad(k\in F).
 \end{split}                                                   \tag{C.81.11}
\]
They are continuous inverse maps in the
specified topology. Translation supplies
the other charts. Thus (C.81.3) is a locally
trivial principal bundle. Its metrizable
manifold base makes it numerable, so the
usual based homotopy lifting and exact
sequence apply to these actual maps.

\emph{The exact positive-degree effect of fixing finitely many orbits.}
For the first orbit the base in (C.81.3)
is the contractible plane $O$. Its homotopy
exact sequence gives an isomorphism on
every positive homotopy group after fixing
that first orbit. Suppose $|Q|\geq2$ and
$A$ is now nonempty. The surface
$M=O\setminus QA$ has at least two punctures.
It is aspherical: a compact sphere image
lies in a compact planar subsurface with
nonempty boundary; that subsurface retracts
onto a finite graph and has no homotopy
in degrees at least two.

Moreover $\operatorname{Homeo}(M)_0$ is
contractible by \cite[Theorem~1.1]{C_Yagasaki2000}.
Its exceptions are explicitly excluded here.
The surface is oriented, and forgetting all
but two punctures gives a surjection onto
a free group of rank two: the two small
meridians and their connecting paths can
be chosen to avoid the remaining locally
finite punctures. Thus it is neither the
plane nor an annulus, and cannot be the
nonorientable exception.

Every loop in $\mathscr C_Q(A)$ restricts
continuously, with inverses, to a loop in
$\operatorname{Homeo}(M)_0$. Its evaluation
at $x$ factors through that restriction,
and hence has zero class in $\pi_1(M,x)$.
Asphericity and the exact sequence of
(C.81.3) now give
\[
 \pi_n(\mathscr C_Q(A\cup\{x\}),1)
       \xrightarrow{\ \cong\ }
 \pi_n(\mathscr C_Q(A),1)
                              \qquad(n\geq1).
                                                               \tag{C.81.12}
\]
For $n=1$, injectivity uses $\pi_2(M)=0$
and surjectivity uses the just proved zero
evaluation homomorphism. The higher degrees
use both adjacent vanishing groups of $M$.
Induction proves the asserted finite-set
comparison for $|Q|\geq2$. For $Q=\{1\}$,
\cite[Lemma~2.1(ii)]{C_Yagasaki2000} applies
directly on the compact disk with its
boundary and the finite fixed set. The
identity components have zero positive
homotopy groups on both sides, giving
the same inclusion isomorphism.

This proves that any sphere class has a
representative fixing any specified finite
number of orbit labels, through an actual
based homotopy in the original relative
group. No uniform choice or limit over
an infinite configuration has been used.

\emph{The explicit comparison with Whitney topology.}
Descent $h\mapsto\bar h$ from $\mathscr C_Q$
is continuous into the compact-open and
inverse topology of $\operatorname{Homeo}(N)$:
local covering sections prove joint continuity
on each compact quotient set, also for
the descended inverse maps.
Assume now that $N$ is noncompact. Choose
points $y_j$ escaping every compact set
and pairwise disjoint, locally finite
coordinate disks $V_j$ around them.
Choose lifts $x_j$. Shrink each disk so
that its lift $U_j$ at $x_j$ has compact
closure, disjoint $Q$-translates, and is
contained in a $d_Q$-ball of radius less
than $1/(4j)$. Then every $qU_j$ has
original $d$-diameter less than $1/j$.

Choose a disk homeomorphism $b_j$ supported
in $V_j$, isotopic to the identity there,
with $b_j(y_j)\ne y_j$. Lift its supported
isotopy from the identity. The endpoint
$h_j$ preserves each lifted disk $qU_j$
and is the identity elsewhere. It has
the actual circle-fixed extension already
proved. Consequently
\[
 \begin{split}
 \sup_{z\in\overline D}d(h_jz,z)&<1/j,\\
 \sup_{z\in\overline D}d(h_j^{-1}z,z)&<1/j.
 \end{split}
                                                               \tag{C.81.13}
\]
Thus $h_j\to1$ in $\mathscr C_Q$ and,
by (C.81.7), in its exact uniform $d_Q$
topology as well.

Put $a_j=\rho(b_j(y_j),y_j)>0$. There is
a continuous positive $\epsilon$ on $N$
with $\epsilon(y_j)<a_j/2$: on each
locally finite disk use a bump function
interpolating between that positive chosen
value at $y_j$ and the value $1$ outside
the disk. These locally finite formulas
give a positive continuous global function.
The one Whitney neighborhood
\[
 \{b:\rho(b(y),y)<\epsilon(y)\text{ for every }y\in N\}
                                                               \tag{C.81.14}
\]
contains none of the $b_j$. Hence the
specified descent homomorphism, even on
the subgroup of actual lifts of compactly
supported quotient isotopies with its
induced $\mathscr C_Q$ topology, is not
continuous to this Whitney topology.

The definition and fragmentation results
in \cite[p.~395 and \S2, Proposition~2.9,
Corollary~2.10]{C_LechMichalikRybicki2018}
use that majorant topology and compactly
supported isotopies. Formula (C.81.13)
and the single excluded neighborhood
(C.81.14) prove the exact comparison
obstruction to their direct use for an
arbitrary original continuous family.
Likewise the metric-cover deformation
in \cite[Theorem~1.1]{C_YagasakiUniform2014}
has a compact quotient hypothesis and
does not itself establish preservation
of our named relative image. The uniform
controlled fragmentation and the coherent
infinite-configuration limit remain the
stated general discrete-case requirements.
\end{proof}
